\documentclass[11pt]{article}

% ---------------------------------------------------------------------------
% WORKSHOP / arXiv preprint draft -- SKELETON ONLY (prose to be authored).
% Swap \documentclass for the venue style file once a target is chosen.
% Story (LOCKED 2026-06-08):
%   Headline   : WE WIN -- the cheapest spectral-polar member is the fastest
%                LoRA optimizer.
%   Protagonist: coupled two-sided DIAGONAL curvature + full polar + op-norm radius
%                (kl-diag-polar-lora).
%   Claim      : it matches the expensive members and beats AdamW 1.2-1.7x
%                across 9 (model,dataset,rank) settings -- the expensive
%                machinery is dead weight.
% ---------------------------------------------------------------------------

\usepackage{amssymb,amsmath,amsthm,amsbsy,amsfonts}
\usepackage{mathtools}
\usepackage[letterpaper,margin=1in]{geometry}
\usepackage{bm}
\usepackage{graphicx}
\usepackage{float}        % [H] keeps trailing appendix figures with their section
\usepackage{booktabs}
\usepackage{algorithm}
\usepackage{algpseudocode}
% Looser line spacing inside algorithm blocks for readability.
\makeatletter
\newcommand{\algspacing}{\setlength{\@tempdima}{0pt}\renewcommand{\baselinestretch}{1.25}\selectfont}
\makeatother
\algrenewcommand\algorithmicrequire{\textbf{Input:}}
% Phase header: blank line + vertical space above, bold label.
\algnewcommand\algorithmicphase{\vspace{5pt}\textbf}
\algnewcommand\Phase{\Statex\algorithmicphase}
\usepackage{enumitem}
\usepackage{xcolor}
\usepackage[colorlinks,citecolor=blue,urlcolor=blue]{hyperref}
\usepackage[nameinlink]{cleveref}
\usepackage[firstinits=true,backend=bibtex,style=numeric-comp,citestyle=numeric-comp,sorting=nyt,url=false,arxiv=false,isbn=false,doi=false]{biblatex}
\bibliography{references}

\graphicspath{{figs/}}

% --- macros -----------------------------------------------------------------
\renewcommand{\vec}[1]{\bm{#1}}
\newcommand{\mat}[1]{\bm{#1}}
\newcommand{\R}{\mathbb{R}}
\newcommand{\E}{\mathbb{E}}
\DeclareMathOperator{\diag}{diag}
\DeclareMathOperator*{\argmin}{arg\,min}
\DeclareMathOperator*{\argmax}{arg\,max}
\newcommand{\smax}{\sigma_{\max}}
\DeclarePairedDelimiter{\norm}{\lVert}{\rVert}
\newcommand{\DW}{\Delta W}
\DeclareMathOperator{\polar}{msign}          % matrix sign (orthogonal polar factor): U\Sigma V^T -> U V^T
% Working method name -- rename in one place once finalized.
\newcommand{\methodname}{\textsc{PoLoRA}\xspace}
\usepackage{xspace}
% Visible marker for numbers pending the full-polar confirm/negate sweep.
\newcommand{\PLACEHOLDER}[1]{\textcolor{red}{[#1]}}
% theorem environments
\theoremstyle{plain}
\newtheorem{lemma}{Lemma}
\newtheorem{proposition}{Proposition}
\newtheorem{corollary}{Corollary}
\theoremstyle{definition}
\newtheorem{definition}{Definition}
\theoremstyle{remark}
\newtheorem{remark}{Remark}

\title{K-FAC Robust Steepest Descent} \author{}
\date{}

\begin{document}
\maketitle

Consider
\[
    \ell(W)=\mathbb{E}_i\,\ell_i(Wx_i).
\]
Define
\[
    b_i
    =
    \left.\nabla_y\ell_i(y)\right|_{y=Wx_i},
    \qquad
    G_i=\nabla_W\ell_i(Wx_i)=b_i x_i^\top,
\]
and let
\[
    G=\mathbb{E}_i G_i,
    \qquad
    g_i=\operatorname{vec}(G_i)=x_i\otimes b_i.
\]
The gradient second moment is
\[
    \Sigma
    =
    \mathbb{E}_i[g_i g_i^\top]
    =
    \mathbb{E}_i
    \left[
        (x_i x_i^\top)\otimes(b_i b_i^\top)
    \right].
\]
K-FAC approximates this matrix by
\[
    \Sigma\approx Q\otimes P,
    \qquad
    Q=\mathbb{E}_i[x_i x_i^\top],
    \qquad
    P=\mathbb{E}_i[b_i b_i^\top].
\]

For a small update $\Delta W$,
\[
    \ell_i((W+\Delta W)x_i)-\ell_i(Wx_i)
    =
    \langle G_i,\Delta W\rangle
    +o(\|\Delta W\|).
\]
Thus $\langle G_i,\Delta W\rangle$ is the first-order change of the
$i$th sample loss. This motivates problem
\[
\begin{array}{ll}
\underset{\Delta W}{\operatorname{minimize}}
    & \langle G,\Delta W\rangle \\
\operatorname{subject\ to}
    & \displaystyle
      \max_i \,\langle G_i,\Delta W\rangle\leq\eta ,
\end{array}
\]
which minimizes the average first-order loss change while limiting the
largest sample-wise increase.

We replace the empirical gradient set
by a rank-one ball in the K-FAC metric. 
For $U=bx^\top$,
\[
\operatorname{vec}(U)^\top
(Q\otimes P)^{-1}
\operatorname{vec}(U)
=
(x\otimes b)^\top
(Q^{-1}\otimes P^{-1})
(x\otimes b) =
(x^\top Q^{-1}x)(b^\top P^{-1}b).
\]
Define
\[
    \mathcal G_\rho
    =
    \left\{
        bx^\top ~\middle|~
        (x^\top Q^{-1}x)(b^\top P^{-1}b)\leq\rho^2
    \right\}.
\]
Choosing
\[
\rho
=
\max_i
\sqrt{
    (x_i^\top Q^{-1}x_i)
    (b_i^\top P^{-1}b_i)
}
\]
ensures that all sample gradients are in the set $\mathcal G_\rho$, therefore,
\[
\max_i \, \langle G_i,\Delta W\rangle
\leq
\max_{U\in\mathcal G_\rho} \,
\langle U,\Delta W\rangle.
\]
Since $\mathcal G_\rho=\rho\mathcal G_1$,
factor $\rho$ can therefore be absorbed into $\eta$. Without
loss of generality, we use $\mathcal G=\mathcal G_1$.


The conservative robust problem is
\[
\begin{array}{ll}
    \underset{\Delta W}{\operatorname{minimize}}
        & \langle G,\Delta W\rangle \\
    \operatorname{subject\ to}
        & \displaystyle
          \max_{U\in\mathcal G} \,
          \langle U,\Delta W\rangle
          \leq\eta .
\end{array}
\]
Let $x=Q^{1/2}u$ and
$b=P^{1/2}v$,
then $U=bx^\top\in\mathcal G$ if and only if
$\|u\|_2\|v\|_2\leq1$, and therefore
\[
    \max_{U\in\mathcal G}\langle U,\Delta W\rangle
    =
    \max_{\|u\|_2\|v\|_2\leq1}
    v^\top P^{1/2}\Delta WQ^{1/2}u 
    =
    \left\|P^{1/2}\Delta WQ^{1/2}\right\|_2.
\]
Finally, the robust steepest descent problem is
\[
\begin{array}{ll}
    \underset{\Delta W}{\operatorname{minimize}}
        & \langle G,\Delta W\rangle \\
    \operatorname{subject\ to}
        &
        \left\|P^{1/2}\Delta WQ^{1/2}\right\|_2
        \leq\eta .
\end{array}
\]


\end{document}
